\documentclass[a4paper,oneside]{article} % Tài liệu dạng sách, in một mặt trên giấy A4
%\documentclass[13pt,a4paper]{extarticle}

% ====== GÓI CƠ BẢN ======
\usepackage{scrextend}
\changefontsizes{13pt}
\usepackage{float,booktabs}
\usepackage[utf8]{vietnam} % Hỗ trợ gõ và hiển thị tiếng Việt
\usepackage[numbers,sort&compress]{natbib} % Gói trích dẫn tài liệu

% ====== DÒNG VĂN BẢN ======
\linespread{1.2} % Giãn dòng
\setlength{\parskip}{\baselineskip} % Khoảng cách giữa các đoạn
\usepackage[skip=6pt, indent=1cm]{parskip} % Tùy chỉnh cách đoạn và thụt đầu dòng
\usepackage{setspace}


% ====== CÀI ĐẶT LỀ TRANG ======
\usepackage[top=2cm,bottom=2cm,right=2cm,left=3cm]{geometry} % Cài đặt lề văn bản và chiều cao phần đầu trang

% ====== GÓI DANH SÁCH ======
\usepackage{enumerate} % Tạo danh sách đánh số
\usepackage{enumitem} % Tuỳ chỉnh kiểu danh sách

% ====== CHÈN HÌNH ẢNH ======
\usepackage{graphicx} % Gói chèn ảnh
\graphicspath{{Figs/}} % Đặt thư mục chứa hình ảnh
\usepackage{pdfpages}
\usepackage{amsmath,amssymb,bm,cool,mathtools,amsfonts,mathptmx} % Các gói công thức Toán

% ====== HYPERLINK ======
\usepackage[hidelinks,unicode,pdfusetitle]{hyperref} % Tạo liên kết nhưng không đổi màu
\hypersetup{
	colorlinks=true,
	citecolor=black,
	linkcolor=black,
	filecolor=black,      
	urlcolor=black,
} % Cấu hình màu liên kết cho in ấn


% ====== THÔNG TIN LUẬN VĂN ======
\newcommand{\university}{ĐẠI HỌC CẦN THƠ}
\newcommand{\school}{Khoa Khoa học Tự nhiên}
\newcommand{\nametitle}{Nhận dạng thống kê cho~hàm~mật~độ~xác~suất mất~cân~bằng và~ứng dụng cho dữ liệu ảnh}
\newcommand{\MainAdvisor}{PGs. TS. Võ Văn Tài}
\newcommand{\student}{Trần Nam Hưng}
\newcommand{\studentid}{M1823002}
\newcommand{\yeardone}{2025}
\newcommand{\major}{Lý~thuyết Xác~suất và Thống~kê Toán~học}
\newcommand{\city}{Cần Thơ}

\begin{document}
\begin{center}
	\begin{tabular}{cc}
		\bf   BỘ GIÁO DỤC VÀ ĐÀO TẠO     &   \bf  CỘNG HÒA XÃ HỘI CHỦ NGHĨA VIỆT NAM\\
 \bf ĐẠI HỌC CẦN THƠ           &     \bf   Độc lập -- Tự do -- Hạnh phúc\\
 \rule{4cm}{1pt} &  \rule{4cm}{1pt}\\
\end{tabular}\\
\vspace{20pt}
{\fontsize{14pt}{12pt}\selectfont\bf BẢN GIẢI TRÌNH CHỈNH SỬA LUẬN VĂN THẠC SĨ}\\
{\it (Thông qua Hội đồng bảo vệ luận văn thạc sĩ ngày 26/10/2025)}\\
\vspace{20pt}
\end{center}
	\begin{tabular*}{\columnwidth}{@{\extracolsep\fill}ll@{\extracolsep\fill}}
	Họ và tên học viên: \selectfont\MakeUppercase\student &                      Mã số học viên: \studentid\\
	Ngành: \major                       &                       Mã ngành: 8 46 01 06\\
\end{tabular*}
	Tên đề tài luận văn Thạc sĩ: \nametitle \\
	Người hướng dẫn: \MainAdvisor, Khoa Khoa học Tự nhiên, Đại học Cần Thơ.
	
	\vspace{20pt}
	Luận văn đã chỉnh sửa theo ý kiến đóng góp của Hội đồng như sau:
	
\begin{table}[H]\centering\footnotesize
	\begin{tabular*}{\columnwidth}{@{\extracolsep\fill}p{0.5cm}p{7.5cm}p{7.5cm}@{\extracolsep\fill}}
		\toprule
		\bf TT &\bf	Nội dung góp ý, đề nghị chỉnh sửa &\bf	Nội dung đã chỉnh sửa\\
		      &\it (Ghi cụ thể nội dung cần chỉnh sửa, mục, trang$\ldots$) &\\
		1 &&\\
		2 &&\\
		\midrule
		
		
	\end{tabular*}
\end{table}
	
\end{document}